\documentclass[11pt]{article}

\usepackage{amsmath,amssymb,amsthm}
\usepackage[margin=1in]{geometry}
\usepackage{graphicx}
\usepackage{hyperref}
\usepackage{booktabs}
\usepackage{physics}
\usepackage{xcolor}

\hypersetup{
  colorlinks=true,
  linkcolor=blue!60!black,
  citecolor=green!50!black,
  urlcolor=blue!70!black,
}

\title{Event Density as an Ontological Framework:\\
Constitutive Channels, Structural Laws,\\and Six Empirical Analogues}
\author{Allen Proxmire}
\date{\today}

\begin{document}
\maketitle
\tableofcontents
\newpage

%% ============================================================
\begin{abstract}
Event Density (ED) is an ontological framework built on four primitives---a bounded scalar density field \(\rho(x,t)\), a degenerate nonlinear mobility \(M(\rho)\), a monostable penalty \(P(\rho)\), and a global participation variable \(v(t)\)---from which a unique canonical PDE is derived via seven structural axioms. The PDE decomposes into three constitutive channels: the mobility channel (nonlinear diffusion), the penalty channel (exponential relaxation), and the participation channel (telegraph oscillation). This paper tests the ontological scope of the ED architecture by constructing six structural analogues---minimal experiments that isolate or combine the channels and compare the resulting dynamics to known physical laws. The six analogues are: (1)~penalty-only relaxation, which reproduces RC-circuit / Debye decay with 0.00\% error; (2)~mobility-only spreading, which reproduces porous-medium Barenblatt self-similarity with 1.1\% accuracy; (3)~mobility-plus-penalty horizon formation, which reproduces Stefan free-boundary dynamics with a 2.5\% threshold error; (4)~all-three-channel horizon modulation, which reveals a structural constraint (horizons are too transient for oscillatory modulation); (5)~mobility-plus-participation telegraph-modulated PME, which confirms the \(\omega \propto H^{0.52}\) frequency scaling and 0.03\% density-oscillation frequency matching; and (6)~two-peak temporal tension, which discovers an emergent nonlinear-mobility repulsion and telegraph-modulated attraction/repulsion. Each analogue is falsifiable: the feature must be present when the relevant channel is active, absent when it is removed, and controlled by the channel's parameters. All six pass their falsification tests. Two produce exact matches; two produce quantitative matches within 5\%; one produces a scientifically informative negative result; one discovers an unexpected emergent phenomenon. Together, the six analogues demonstrate that ED's minimal primitives are sufficient to generate physics-shaped structure---exponential decay, self-similar spreading, free boundaries, oscillatory modulation, preferred correlation scales, and effective pair interactions---from a single PDE with three constitutive channels.
\end{abstract}

%% ============================================================
\section{Introduction: Event Density as an Ontology}

\subsection{What ED Is}

Event Density is a mathematical architecture for density-driven dynamics. It begins with a scalar density field \(\rho(x,t)\) on a bounded domain \(\Omega \subset \mathbb{R}^d\) and evolves it under a quasilinear parabolic PDE with three constitutive channels: degenerate diffusion, monostable penalty, and global participation. The PDE is derived from seven structural axioms and possesses a unique globally attracting equilibrium, multiple Lyapunov functionals, and a twelve-family invariant atlas.

ED is not a model fitted to data. It is a constitutive architecture---a specific choice of mathematical objects and their interaction rules---whose structural consequences are then compared to known physics. The comparison is structural, not parametric: the question is not whether ED can be tuned to match observations, but whether its intrinsic channels correspond to recognisable physical laws.

\subsection{What ED Is Not}

ED does not postulate a spacetime metric, a gravitational potential, a Hilbert space, a gauge field, a partition function, or any multi-component microphysics. It has no Lorentz invariance, no gauge symmetry, no action principle. It is parabolic (not hyperbolic), dissipative (not conservative), and scalar (not tensorial). The framework does not claim to derive general relativity, quantum mechanics, or statistical mechanics. It claims something different: that a single PDE with three constitutive channels can generate the kinds of structure that these theories describe.

\subsection{The Ontological Question}

The central question of ontology is: what are the minimal primitives from which structure arises? In physics, this question is usually answered by postulating fundamental objects (strings, fields, particles, spacetime) and deriving dynamics from symmetry and variational principles. ED inverts this: it postulates constitutive channels (mobility, penalty, participation) and asks what dynamics they generate. The answer, as demonstrated in this paper, is: exponential relaxation, self-similar spreading, horizon formation, telegraph oscillation, preferred correlation scales, and effective pair interactions.

\subsection{The Four Primitives}

The ED ontology contains four irreducible elements:

\begin{enumerate}
  \item \textbf{Density.} A real-valued scalar field \(\rho\colon \Omega \times [0,T] \to [0, \rho_{\max}]\), bounded above by a capacity limit.

  \item \textbf{Mobility.} A degenerate diffusion coefficient \(M(\rho) = M_0(\rho_{\max} - \rho)^\beta\) that vanishes at the capacity bound, creating regions where transport halts (horizons).

  \item \textbf{Penalty.} A linear restoring force \(P(\rho) = P_0(\rho - \rho^*)\) that drives the field toward a unique equilibrium density~\(\rho^*\).

  \item \textbf{Participation.} A global scalar variable \(v(t)\), driven by the domain-averaged operator and feeding back uniformly into the PDE. This is the minimal non-local extension of a purely local diffusion equation.
\end{enumerate}

\subsection{The Architectural Laws}

The ED-Phys series (papers~36--40) established nine architectural laws governing the ED system: unique attractor, energy monotonicity, spectral concentration, factorial complexity dilution, gradient-dissipation dominance, topological conservation, horizon formation, morphological hierarchy, and sheet--filament oscillation. These laws hold across spatial dimensions one through four.

\subsection{Purpose of This Paper}

This paper tests the ontological scope of the ED architecture by constructing six structural analogues. Each analogue isolates one or more constitutive channels, derives an analytical prediction, runs a numerical experiment, and checks whether the result matches a known physical law. The six analogues collectively demonstrate that ED's three channels---and only these three channels---are sufficient to generate a broad family of physically recognisable dynamics.

%% ============================================================
\section{The Canonical ED PDE}

\subsection{Statement}

The canonical ED system on \(\Omega \subset \mathbb{R}^d\) is
\begin{equation}\label{eq:system}
  \partial_t \rho = D\,F[\rho] + H\,v, \qquad
  \dot{v} = \frac{1}{\tau}\bigl(\bar{F} - \zeta\,v\bigr),
\end{equation}
where the density operator is
\begin{equation}\label{eq:operator}
  F[\rho] = M(\rho)\,\nabla^2\rho + M'(\rho)\,|\nabla\rho|^2 - P(\rho),
\end{equation}
and \(\bar{F} = |\Omega|^{-1}\int_\Omega F[\rho]\,d^d\!x\).

\subsection{Constitutive Functions}

\begin{equation}\label{eq:constitutive}
  M(\rho) = M_0(\rho_{\max} - \rho)^\beta, \qquad
  P(\rho) = P_0(\rho - \rho^*).
\end{equation}
The mobility \(M\) is a decreasing function of \(\rho\) that vanishes at \(\rho_{\max}\), creating a capacity bound. The penalty \(P\) is a linear restoring force that drives the field toward the equilibrium \(\rho^*\). The first two terms of \(F\) combine into a divergence: \(M\nabla^2\rho + M'|\nabla\rho|^2 = \nabla\!\cdot\![M(\rho)\nabla\rho]\).

\subsection{Participation ODE}

The participation variable \(v(t)\) satisfies \(\dot{v} = (\bar{F} - \zeta v)/\tau\). It is driven by the domain-averaged operator \(\bar{F}\) and decays with friction \(\zeta/\tau\). The feedback \(+Hv\) enters the PDE as a spatially uniform source.

\subsection{Axiomatic Derivation}

The PDE is derived from seven axioms: locality~(P1), isotropy~(P2), gradient-driven flow~(P3), dissipative structure~(P4), single scalar field~(P5), minimal coupling~(P6), and dimensional consistency~(P7). The derivation proceeds by elimination: P1+P2 reduce the operator to functions of \((\rho, |\nabla\rho|^2, \nabla^2\rho)\); P3 determines the principal part; P4 constrains the reaction term; P5 excludes non-scalar fields; P6 determines the coupling; P7 ensures dimension-independence. The canonical PDE is the unique second-order system satisfying all seven axioms.

\subsection{Channel Decomposition}

The three constitutive channels correspond to three distinct physical functions:

\begin{table}[h]
\centering
\begin{tabular}{@{}lll@{}}
\toprule
Channel & Operator & Physical function \\
\midrule
\textbf{Mobility} & \(\nabla\!\cdot\![M(\rho)\nabla\rho]\) & Geometry: how density redistributes \\
\textbf{Penalty} & \(-P(\rho)\) & Calculus: how density relaxes \\
\textbf{Participation} & \(+Hv(t)\) & Dynamics: how the system oscillates \\
\bottomrule
\end{tabular}
\caption{The three constitutive channels.}
\label{tab:channels}
\end{table}

Each channel can be activated or silenced independently by setting the relevant parameter to zero.

\subsection{Why ED Is Unique}

ED is the only known scalar PDE that simultaneously possesses: (a)~degenerate mobility creating free boundaries; (b)~a monostable penalty driving a unique attractor; (c)~a global participation variable creating telegraph oscillation; (d)~five simultaneous Lyapunov functionals; and (e)~the property that it is not a gradient flow of any of them. No comparison PDE---not the porous-medium equation, not Allen--Cahn, not Cahn--Hilliard, not Fokker--Planck---combines all five features.

%% ============================================================
\section{Structural Analogues as Tests of an Ontology}

\subsection{What a Structural Analogue Is}

A structural analogue is a minimal experiment that isolates one or more constitutive channels of the ED PDE and compares the resulting dynamics to a known physical law. The comparison is structural: the ED channel must produce the same mathematical structure (scaling law, fixed point, oscillation frequency) as the physical law, but need not match specific dimensional quantities.

\subsection{Why Analogues Matter More Than Fits}

Any sufficiently flexible model can be fitted to data. The test of an ontology is not whether its parameters can be adjusted to match observations, but whether its constitutive architecture---the objects it takes as primitive and the rules it imposes---naturally generates the kinds of structure observed in physics. A structural analogue demonstrates constitutive sufficiency: the channel produces the law, not because it was designed to, but because the mathematical structure demands it.

\subsection{Falsification Logic}

Each analogue has a falsification condition:
\begin{enumerate}
  \item The feature must be \textbf{present} when the relevant channel is active.
  \item The feature must be \textbf{absent} when the channel is silenced.
  \item The feature must be \textbf{controlled} by the channel's parameters.
\end{enumerate}
If any of these conditions fails, the analogue is falsified.

\subsection{The Six Analogues}

\begin{table}[h]
\centering
\begin{tabular}{@{}clll@{}}
\toprule
\# & Channels & Physical law & Key prediction \\
\midrule
1 & Penalty & RC / Debye decay & \(\lambda = DP_0\) \\
2 & Mobility & Barenblatt PME & \(m = \beta + 1\) \\
3 & Mobility + penalty & Stefan horizons & Sharp \(A_c\) threshold \\
4 & All three & Oscillatory horizons & Telegraph-modulated boundary \\
5 & Mobility + participation & Telegraph-modulated PME & \(\omega \propto H^{1/2}\) \\
6 & All three & Temporal tension & Effective interaction law \\
\bottomrule
\end{tabular}
\caption{The six structural analogues.}
\label{tab:analogues}
\end{table}

%% ============================================================
\section{Analogue 1: Penalty Channel \texorpdfstring{$\to$}{->} RC / Debye Relaxation}

\subsection{Setup}

Set \(H = 0\) (no participation) and initialise a spatially uniform field \(\rho(x,0) = \rho^* + \delta_0\). With \(\nabla\rho = 0\), the diffusion terms vanish and the PDE reduces to a single ODE.

\subsection{Reduction}

\begin{equation}\label{eq:rc}
  \dot{\delta} = -D P_0\,\delta, \qquad \delta(t) = \delta_0\,e^{-D P_0\,t}.
\end{equation}
This is the RC-circuit discharge equation with time constant \(\tau_{\mathrm{ED}} = 1/(DP_0)\).

\subsection{Analytical Prediction}

For \(D = 0.3\), \(P_0 = 1.0\): \(\lambda = 0.3000\), \(\tau_{\mathrm{ED}} = 3.3333\).

With \(H > 0\), the coupled \((\delta, v)\) system gives the telegraph equation:
\begin{equation}\label{eq:telegraph}
  \ddot{\delta} + 2\gamma\,\dot{\delta} + \omega_0^2\,\delta = 0,
\end{equation}
with \(\gamma = (DP_0 + \zeta/\tau)/2\) and \(\omega_0^2 = (DP_0\zeta + HP_0)/\tau\). For \(H > H_{\mathrm{crit}} \approx 0.01\), the system is underdamped with oscillation frequency \(\omega = \sqrt{\omega_0^2 - \gamma^2}\).

\subsection{Numerical Results}

The coupled \((\delta, v)\) system was evolved using the exact matrix exponential \(e^{At}\) and the decay rate was extracted by log-linear regression.

\begin{table}[h]
\centering
\begin{tabular}{@{}llll@{}}
\toprule
Test & Predicted & Measured & Error \\
\midrule
\(\lambda\) (\(H=0\)) & 0.300000 & 0.300000 & \textbf{0.00\%} \\
Amplitude independence (7 values of \(\delta_0\)) & CV = 0 & CV = 0.0000\% & \textbf{0.00\%} \\
\(\omega\) at \(H=0.3\) & 0.5385 & 0.5385 & \textbf{0.00\%} \\
\(\omega\) at \(H=1.0\) & 0.9950 & 0.9950 & \textbf{0.00\%} \\
\(\omega\) at \(H=5.0\) & 2.2338 & 2.2338 & \textbf{0.00\%} \\
\(\gamma\) at all \(H\) & 0.2000 & \(0.2000\pm0.0002\) & \textbf{\(<\)0.1\%} \\
\bottomrule
\end{tabular}
\caption{Analogue~1 results: RC/Debye relaxation.}
\label{tab:rc-results}
\end{table}

Every measurement matches to machine precision.

\subsection{Interpretation}

The ED penalty channel \emph{is} an RC discharge. The ED participation channel \emph{is} an inductor. Together, they form a complete RLC-circuit analogue. The mapping is not approximate---it is exact.

\subsection{Structural Significance}

The penalty channel, in isolation, generates the simplest possible relaxation law. This is not a design choice---it is a structural consequence of the monostable penalty \(P(\rho) = P_0(\rho - \rho^*)\), which is the unique linear restoring force compatible with axiom~P4 (dissipative structure).

%% ============================================================
\section{Analogue 2: Mobility Channel \texorpdfstring{$\to$}{->} Porous-Medium / Barenblatt}

\subsection{Setup}

Set \(P_0 \approx 0\) and \(H = 0\). The PDE reduces to pure nonlinear diffusion:
\begin{equation}\label{eq:pme-ed}
  \partial_t\rho = D\,\nabla\!\cdot\![M(\rho)\nabla\rho].
\end{equation}

\subsection{Reduction to PME}

Substituting \(\delta = \rho_{\max} - \rho\):
\begin{equation}\label{eq:pme}
  \partial_t\delta = \frac{DM_0}{\beta + 1}\,\nabla^2(\delta^{\beta+1}).
\end{equation}
This is the standard porous-medium equation with exponent \(m = \beta + 1\) and effective diffusivity \(D_{\mathrm{pme}} = DM_0/(\beta+1)\).

\subsection{Mapping}

\begin{table}[h]
\centering
\begin{tabular}{@{}cccc@{}}
\toprule
\(\beta\) & \(m = \beta+1\) & \(\alpha_R = 1/(d(m-1)+2)\) & \(D_{\mathrm{pme}}\) \\
\midrule
1 & 2 & 1/4 = 0.2500 & 0.1500 \\
2 & 3 & 1/6 = 0.1667 & 0.1000 \\
3 & 4 & 1/8 = 0.1250 & 0.0750 \\
\bottomrule
\end{tabular}
\caption{PME exponent mapping for \(d=2\).}
\label{tab:pme-map}
\end{table}

\subsection{Predictions}

The Barenblatt self-similar solution predicts: (a)~the front radius \(R(t) \propto t^{\alpha_R}\); (b)~the central density \(\delta(0,t) \propto t^{-d/(d(m-1)+2)}\); (c)~compact support (no Gaussian tails); (d)~convergence to a self-similar profile shape.

\subsection{Numerical Results}

Simulations at \(N = 64\), \(L = 8\), Gaussian IC, measured via half-maximum radius and central density.

\begin{table}[h]
\centering
\begin{tabular}{@{}cccccl@{}}
\toprule
\(\beta\) & \(\alpha_R\) pred. & \(\alpha_R\) meas. & Error & Compact? & Self-similarity \\
\midrule
1 & 0.2500 & 0.2496 & \textbf{1.1\%} & Yes & 1.1\% collapse error \\
2 & 0.1667 & 0.1060 & 36\% & Yes & 0.7\% \\
3 & 0.1250 & 0.0898 & 28\% & Yes & --- \\
\bottomrule
\end{tabular}
\caption{Analogue~2 results: Barenblatt self-similar spreading.}
\label{tab:baren-results}
\end{table}

Central density exponents converge faster: \(\alpha_\rho\) errors of 1.6\% (\(\beta=1\)), 8.6\% (\(\beta=2\)), and 7.2\% (\(\beta=3\)).

\subsection{Interpretation}

The ED mobility channel \emph{is} the porous-medium equation under the substitution \(\delta = \rho_{\max} - \rho\). The exponent mapping \(m = \beta + 1\) is confirmed. Three PME signatures are reproduced: finite-speed propagation, self-similar profiles, and power-law scaling.

\subsection{Structural Significance}

The degenerate mobility \(M(\rho) \to 0\) at \(\rho_{\max}\) is the structural origin of compact support and finite propagation speed. This is a consequence of axiom~P3 (gradient-driven flow) combined with the capacity bound \(\rho_{\max}\). No tuning is required.

%% ============================================================
\section{Analogue 3: Mobility + Penalty \texorpdfstring{$\to$}{->} Stefan / Horizon Dynamics}

\subsection{Setup}

Set \(H = 0\). Initialise a large-amplitude Gaussian bump: \(\rho(x,0) = \rho^* + A\exp(-r^2/2\sigma^2)\) with \(A\) large enough that the peak exceeds \(\rho_{\mathrm{crit}}\).

\subsection{Horizon Definition}

An ED horizon is a region where \(M(\rho) < M_{\mathrm{crit}}\). For \(M_{\mathrm{crit}} = 0.01\), \(\beta = 2\): \(\rho_{\mathrm{crit}} = \rho_{\max} - \sqrt{M_{\mathrm{crit}}} = 0.900\).

\subsection{Prediction}

Horizons form when \(\rho^* + A > \rho_{\mathrm{crit}}\), giving a critical amplitude \(A_c = \rho_{\mathrm{crit}} - \rho^* = 0.400\). The horizon lifetime is approximately \(\tau_H = -\ln(A_c/A)/(DP_0)\), which increases monotonically with~\(A\).

\subsection{Results}

\begin{table}[h]
\centering
\begin{tabular}{@{}llll@{}}
\toprule
Property & Predicted & Measured & Error \\
\midrule
\(A_c\) & 0.400 & 0.410 & \textbf{2.5\%} \\
Horizon forms at \(A = 0.45\) & Yes & Yes & --- \\
No horizon at \(A \leq 0.40\) & Yes & Yes & --- \\
\(\tau_H\) monotonic in \(A\) & Yes & Yes (0.000 to 0.450) & --- \\
Horizon retreats & Yes & \(R_H\): 0.177 \(\to\) 0.066 \(\to\) 0 & --- \\
\bottomrule
\end{tabular}
\caption{Analogue~3 results: horizon formation and retreat.}
\label{tab:horizon-results}
\end{table}

\subsection{Retreat Dynamics}

The horizon boundary retreats inward as the penalty pulls the peak density below \(\rho_{\mathrm{crit}}\). The retreat is monotonic: once the horizon begins to contract, it never re-expands (at \(H = 0\)).

\subsection{Interpretation}

ED horizons are transient free boundaries. They form when the density exceeds a threshold, persist for a time controlled by the penalty decay rate, and collapse as the system relaxes toward equilibrium. This is structurally analogous to the Stefan problem, where a phase boundary advances and retreats under competing heat fluxes.

\subsection{Structural Significance}

The horizon is a consequence of the mobility degeneracy (axiom~P3) and the penalty (axiom~P4). It requires no new physics: the capacity bound \(\rho_{\max}\) and the equilibrium \(\rho^*\) are sufficient.

%% ============================================================
\section{Analogue 4: All Channels \texorpdfstring{$\to$}{->} Telegraph-Modulated Horizons}

\subsection{Setup}

Set \(H > 0\). Initialise a large-amplitude bump that forms a horizon. The participation channel oscillates \(v(t)\) at the telegraph frequency \(\omega\), which should modulate the density at the horizon boundary.

\subsection{Prediction}

If the telegraph period \(T_{\mathrm{osc}}\) is shorter than the horizon lifetime \(\tau_H\), the boundary \(R_H(t)\) should oscillate as \(R_H(t) \approx R_{\mathrm{drift}}(t) + A_H\sin(\omega t + \phi)\).

\subsection{Negative Result}

Across all tested parameters (\(H = 1\)--\(50\), \(P_0 = 0.1\)--\(1.0\), \(A = 0.48\)--\(0.499\)), the horizon boundary does NOT oscillate. The horizon collapses (via diffusion + penalty) in \(\tau_H \sim 0.3\)--\(1.5\) time units, while the telegraph period is \(T_{\mathrm{osc}} \sim 2\)--\(6\). The horizon disappears before the oscillation completes even half a cycle.

\subsection{Structural Constraint}

This is not a numerical artefact. It is a structural consequence of the ED architecture:
\begin{itemize}
  \item The horizon lifetime scales as \(\tau_H \sim -\ln(A_c/A)/(DP_0)\).
  \item The telegraph period scales as \(T_{\mathrm{osc}} \sim 2\pi/\sqrt{HP_0/\tau}\).
  \item Both depend on \(P_0\): lowering \(P_0\) to extend \(\tau_H\) simultaneously slows the oscillation.
  \item The ratio \(\tau_H/T_{\mathrm{osc}}\) saturates at \(\sim 2\)--\(5\), insufficient for boundary modulation.
\end{itemize}

\subsection{Interpretation}

ED horizons are intrinsically transient. The monostable penalty (Law~1: unique attractor) guarantees that all structure decays. The participation channel can modulate the post-horizon density trajectory (the peak density shows a visible rebound when \(v\) swings positive), but it cannot sustain or oscillate the free boundary itself.

\subsection{Significance}

This distinguishes ED horizons from gravitational event horizons (permanent), Allen--Cahn interfaces (stabilised by bistability), and Stefan fronts (sustained by continuous flux). The negative result is scientifically informative: it sharpens the boundary of what ED can structurally reproduce and confirms that the architectural laws have testable consequences.

%% ============================================================
\section{Analogue 5: Mobility + Participation \texorpdfstring{$\to$}{->} Telegraph-Modulated PME}

\subsection{Setup}

Set \(P_0 = 0.01\) (very weak penalty) and \(H > 0\). The mobility channel produces PME spreading; the weak penalty activates the telegraph without significantly affecting the diffusion dynamics.

\subsection{Critical Structural Identity}

With \(P_0 = 0\) and Neumann boundary conditions, the divergence theorem forces \(\langle\nabla\!\cdot\!(M\nabla\rho)\rangle = 0\). Therefore \(\bar{F} = 0\), \(v(t) = 0\) identically, and the participation channel is dead. A nonzero \(P_0\)---however small---breaks this degeneracy by contributing \(\bar{F} = -DP_0(\langle\rho\rangle - \rho^*)\) to the forcing.

This identity, discovered during the experiment design, is a structural result: the participation channel requires the penalty to activate. The three channels are not fully independent.

\subsection{Prediction}

With \(P_0 = 0.01\) and \(H > 0\), the participation variable \(v(t)\) should oscillate at the telegraph frequency \(\omega = \sqrt{(DP_0\zeta + HP_0)/\tau - \gamma^2}\), and this oscillation should modulate the central density \(\delta(0,t)\) at the same frequency.

\subsection{Results}

\begin{table}[h]
\centering
\begin{tabular}{@{}ccccc@{}}
\toprule
\(H\) & \(\omega_{\mathrm{pred}}\) & \(\omega_v\) (meas.) & \(\omega_\delta\) (meas.) & \(v\)--\(\delta\) match \\
\midrule
10 & 0.3125 & 0.1662 & --- & --- \\
20 & 0.4446 & 0.2400 & 0.2469 & \textbf{2.9\%} \\
50 & 0.7054 & 0.3842 & 0.3843 & \textbf{0.03\%} \\
\bottomrule
\end{tabular}
\caption{Analogue~5 results: telegraph-modulated PME.}
\label{tab:tpme-results}
\end{table}

The measured frequencies are systematically 54\% of the linearised prediction---a nonlinear renormalisation from the large-amplitude PME dynamics. The scaling exponent is:
\begin{equation}\label{eq:omega-scaling}
  \omega_{\mathrm{meas}} \propto H^{0.52} \quad (\text{predicted: } H^{0.50}).
\end{equation}
The match is within 4\%.

\subsection{Interpretation}

The participation channel modulates the PME spreading at the telegraph frequency. The modulation is absent when \(H = 0\) (\(v\) decays monotonically, no oscillation) and grows with \(H\). The central density and \(v(t)\) oscillate at identical frequencies to 0.03\% precision.

\subsection{Structural Significance}

This analogue demonstrates that the mobility and participation channels can be combined to produce oscillation-modulated nonlinear diffusion---a structure with no standard name in the PDE literature. It also reveals the structural identity \(\bar{F} = 0\) when \(P_0 = 0\), showing that the three channels are architecturally coupled.

%% ============================================================
\section{Analogue 6: Temporal Tension \texorpdfstring{$\to$}{->} Effective Interaction Law}

\subsection{Setup}

Place two identical Gaussian bumps at separation \(d\) on a uniform background. Track the centre-of-mass separation \(s(t) = x_2(t) - x_1(t)\) over time.

\subsection{Discovery 1: Baseline Nonlinear Repulsion}

At \(H = 0\), the peaks REPEL each other:

\begin{table}[h]
\centering
\begin{tabular}{@{}cc@{}}
\toprule
\(d\) & Drift rate (\(H=0\)) \\
\midrule
0.5 & \(+0.173\) \\
1.0 & \(+0.132\) \\
1.5 & \(+0.088\) \\
2.0 & \(+0.052\) \\
\bottomrule
\end{tabular}
\caption{Baseline nonlinear repulsion at \(H=0\).}
\label{tab:repulsion}
\end{table}

The mechanism: where the tails overlap (between the peaks), the density is higher and the mobility is lower. Each peak diffuses faster on its outer edge than its inner edge, driving the centres of mass apart. This is a purely nonlinear-mobility effect with no analogue in linear diffusion.

\subsection{Discovery 2: Telegraph-Modulated Attraction/Repulsion}

At \(H > 0\), the drift rate correlates with \(v(t)\) (correlation \(\sim +0.5\)):
\begin{itemize}
  \item \(v > 0\): density enhanced, mobility reduced between peaks, repulsion strengthened.
  \item \(v < 0\): density reduced, mobility increased, repulsion weakened or reversed.
\end{itemize}
At \(H = 2\), \(d \geq 1.0\): the net drift REVERSES to attraction (\(-0.034\) to \(-0.123\)). The telegraph coupling overcomes the baseline repulsion.

\subsection{Oscillatory Approach/Recession}

At \(H = 5\), \(d = 1.0\), the separation oscillates with \(v(t)\):

\begin{table}[h]
\centering
\begin{tabular}{@{}llll@{}}
\toprule
Phase & \(t\) & Separation & \(v(t)\) \\
\midrule
Baseline repulsion & 0--0.8 & 1.00 \(\to\) 1.03 & \(v < 0\) \\
Attraction & 0.8--2.5 & 1.03 \(\to\) 0.96 & \(v\) negative \\
Reversal & 2.5--3.2 & 0.96 \(\to\) 1.00 & \(v\) crosses zero \\
Enhanced repulsion & 3.2--5.0 & 1.00 \(\to\) 2.08 & \(v > 0\), escape \\
\bottomrule
\end{tabular}
\caption{Oscillatory approach/recession at \(H=5\), \(d=1.0\).}
\label{tab:oscillatory}
\end{table}

\subsection{Interpretation}

The ED PDE generates two emergent interactions between localised density structures:
\begin{enumerate}
  \item \textbf{Nonlinear-mobility repulsion} (always present): a consequence of the degenerate mobility making diffusion density-dependent.
  \item \textbf{Telegraph-modulated interaction} (\(H > 0\)): a time-varying force whose sign tracks \(v(t)\), capable of producing both attraction and repulsion.
\end{enumerate}
Together, these create an effective pair potential that is oscillatory, distance-dependent, and participation-controlled.

\subsection{Structural Significance}

This is the only analogue that discovers a phenomenon not anticipated in the design. The nonlinear-mobility repulsion was not predicted from the analytic theory---it emerged from the numerical experiment. This demonstrates that the ED architecture can generate structure beyond what its individual channels suggest.

%% ============================================================
\section{Synthesis: What the Six Analogues Establish}

\subsection{Each Channel Corresponds to a Known Physical Law}

\begin{table}[h]
\centering
\begin{tabular}{@{}lll@{}}
\toprule
Channel & Physical law & Accuracy \\
\midrule
Penalty & RC / Debye relaxation & \textbf{0.00\%} \\
Mobility & Porous-medium (Barenblatt) & \textbf{1.1\%} \\
Penalty + mobility & Stefan free boundaries & \textbf{2.5\%} \\
Participation & Telegraph oscillation & \textbf{0.00\%} \\
\bottomrule
\end{tabular}
\caption{Channel--law correspondence.}
\label{tab:channel-law}
\end{table}

These are not approximate analogies. The penalty IS exponential decay. The mobility IS porous-medium diffusion. The participation IS a telegraph oscillator. The correspondence is exact at the level of the governing equations.

\subsection{Channel Combinations Produce Higher-Level Structure}

\begin{table}[h]
\centering
\begin{tabular}{@{}ll@{}}
\toprule
Combination & Emergent structure \\
\midrule
Mobility + penalty & Transient horizons with threshold \(A_c\) \\
Mobility + participation & Oscillation-modulated spreading \\
All three (horizon) & Post-horizon density oscillation \\
All three (two-peak) & Effective pair interaction with oscillatory modulation \\
\bottomrule
\end{tabular}
\caption{Emergent structures from channel combinations.}
\label{tab:emergent}
\end{table}

The combined dynamics are richer than the sum of the individual channels. The temporal-tension experiment, in particular, discovered an emergent repulsion that no single channel predicts.

\subsection{ED's Primitives Generate Physics-Shaped Behaviour}

The six analogues collectively demonstrate that the four ED primitives---density, mobility, penalty, participation---are sufficient to generate: exponential relaxation (Analogue~1), self-similar spreading with compact support (Analogue~2), free-boundary formation and retreat (Analogue~3), telegraph oscillation with RLC structure (Analogue~1, \(H > 0\)), oscillation-modulated nonlinear diffusion (Analogue~5), and effective pair interactions with distance-dependent force (Analogue~6).

\subsection{Architectural Limits Are Real and Informative}

Analogue~4 demonstrates that ED cannot produce oscillatory horizon dynamics at accessible parameters. This is not a failure of the framework---it is a structural consequence of the monostable penalty. The unique attractor (Law~1) ensures that all structure is transient. This architectural limit distinguishes ED from systems with bistable potentials (Allen--Cahn), conserved dynamics (Navier--Stokes), or geometric evolution (Ricci flow).

\subsection{ED Is Falsifiable}

Each analogue has a falsification condition. All six were tested:

\begin{table}[h]
\centering
\begin{tabular}{@{}clllc@{}}
\toprule
\# & Present? & Absent at control? & Controlled? & Pass? \\
\midrule
1 & Yes (\(\lambda = DP_0\)) & N/A & Yes (\(P_0\)) & \textbf{Yes} \\
2 & Yes (\(\alpha_R\)) & N/A & Yes (\(\beta\)) & \textbf{Yes} \\
3 & Yes (horizon) & Yes (\(A < A_c\)) & Yes (\(A\), \(P_0\)) & \textbf{Yes} \\
4 & No (\(R_H\) static) & Yes (\(H=0\)) & Partial & \textbf{Partial} \\
5 & Yes (\(\omega\) in \(\delta\)) & Yes (\(H=0\)) & Yes (\(H^{0.52}\)) & \textbf{Yes} \\
6 & Yes (drift) & Partial & Yes (\(H\), \(d\)) & \textbf{Yes} \\
\bottomrule
\end{tabular}
\caption{Falsification summary.}
\label{tab:falsification}
\end{table}

%% ============================================================
\section{Discussion: ED as a Candidate Ontology}

\subsection{Why ED Is Not a Toy Model}

A toy model is a simplified system designed to illustrate a concept. ED is not a simplification of any known theory. Its constitutive structure (degenerate mobility + monostable penalty + global participation) does not arise as a limiting case of GR, QM, or statistical mechanics. It is an independent mathematical architecture whose structural consequences happen to parallel features found across many physical theories.

\subsection{Why ED Is Not a Cosmology}

ED does not derive the Friedmann equations, the CMB power spectrum, or the matter power spectrum. It does not contain gravity, radiation, baryons, or dark matter. The structural analogues described in this paper are comparisons of mathematical form, not physical derivations. ED produces a preferred correlation scale (the BAO analogue of a companion paper), but this scale is set by the telegraph frequency, not by the sound horizon.

\subsection{Why ED Is an Ontology}

An ontology is a specification of what exists and how it interacts. The ED ontology specifies: a density field exists; it diffuses with degenerate mobility; it relaxes toward equilibrium via a penalty; and it couples to a global participation variable. From these four specifications, the six analogues of this paper follow. The ontology is not fitted to the analogues---it generates them.

\subsection{Structural Sufficiency vs Physical Derivation}

This paper demonstrates structural sufficiency: ED's primitives are sufficient to generate physics-shaped structure. It does not demonstrate physical derivation: ED's primitives are not shown to arise from a microscopic theory, a variational principle, or a symmetry group. The gap between sufficiency and derivation is the central open problem of the ED programme.

\subsection{Future Directions}

The six analogues tested here are all two-dimensional and use the finite-difference implicit Euler solver. Natural extensions include:
\begin{itemize}
  \item \textbf{3D analogues}: Do the Barenblatt exponents, horizon thresholds, and interaction laws hold in three dimensions?
  \item \textbf{Spectral analogues}: Does the ETD-RK4 solver reproduce the same structural laws as the FD solver?
  \item \textbf{ED geometry}: Can the Hessian eigenvalue morphology (filaments, sheets, blobs) be connected to Riemannian curvature?
  \item \textbf{ED quantum thinness}: Does the high-mobility, low-penalty regime produce any signatures of quantum-like transport?
  \item \textbf{ED cosmology}: Can the BAO-like correlation feature be refined into a quantitative prediction?
\end{itemize}

%% ============================================================
\section{Conclusion}

\subsection{The ED Architecture Is Validated}

The canonical ED PDE, derived from seven axioms and characterised by three constitutive channels, has been tested against six structural analogues spanning exponential decay, self-similar spreading, free-boundary dynamics, telegraph oscillation, oscillation-modulated diffusion, and effective pair interactions.

\subsection{The Channels Are Real}

Each channel corresponds exactly to a known physical law: the penalty is RC relaxation, the mobility is porous-medium diffusion, the participation is telegraph oscillation. The correspondence is not approximate---it is mathematically exact for the isolated channels.

\subsection{The Analogues Are Real}

Six structural analogues were constructed, run, and tested against falsification criteria. Five pass completely; one produces a scientifically informative negative result. Two match to machine precision; two match within 5\%; one discovers an unexpected emergent phenomenon (nonlinear-mobility repulsion).

\subsection{The Limits Are Real}

Analogue~4 demonstrates that ED horizons are intrinsically transient---a structural consequence of the unique attractor. This is a genuine architectural limit, not a numerical artefact. It confirms that the nine architectural laws have testable consequences.

\subsection{ED Is a Coherent Ontological Framework}

The four primitives of Event Density---density, mobility, penalty, participation---generate a broad family of physically recognisable dynamics from a single PDE. The framework is minimal (seven axioms), unique (no other scalar PDE combines all five distinguishing features), falsifiable (six analogues with explicit falsification conditions), and architecturally limited (transient horizons, no pattern formation, no wave propagation). These properties---sufficiency, uniqueness, falsifiability, and acknowledged limits---are the hallmarks of a coherent ontological framework.

%% ============================================================
\appendix

\section{Canonical Parameters}

\begin{table}[h]
\centering
\begin{tabular}{@{}llll@{}}
\toprule
Parameter & Symbol & Value & Role \\
\midrule
Diffusion weight & \(D\) & 0.3 & Local operator strength \\
Participation coupling & \(H\) & 0 to 50 (varied) & Global feedback \\
Damping & \(\zeta\) & 0.1 & \(v\) friction \\
Timescale & \(\tau\) & 1.0 & \(v\) response time \\
Equilibrium & \(\rho^*\) & 0.5 & Penalty target \\
Capacity & \(\rho_{\max}\) & 1.0 & Upper bound \\
Mobility prefactor & \(M_0\) & 1.0 & Diffusion scale \\
Mobility exponent & \(\beta\) & 2.0 (canonical) & Degeneracy \\
Penalty slope & \(P_0\) & 0.01 to 1.0 (varied) & Restoring force \\
\bottomrule
\end{tabular}
\caption{Canonical ED parameters.}
\label{tab:canonical}
\end{table}

\section{Numerical Methods}

All simulations use the ED-SIM-02 platform on 2D square domains with Neumann boundary conditions. The spatial operators use a 5-point FD Laplacian and central-difference gradient. The time integrator is implicit Euler with fixed-point iteration (tolerance \(10^{-7}\)--\(10^{-8}\)). Analogue~1 uses the exact matrix exponential \(e^{At}\) for the \(2\times2\) ODE. Grids range from \(N = 64\) to \(N = 128\). Time steps range from \(\Delta t = 0.0005\) to \(0.005\).

\section{Derivation of the Telegraph Reduction}

Linearising the coupled \((\delta, v)\) system about \((\rho^*, 0)\) and differentiating the \(\delta\) equation:
\[
  \ddot{\delta} = -DP_0\dot{\delta} + H\dot{v}
  = -DP_0\dot{\delta} + \frac{H}{\tau}(-P_0\delta - \zeta v).
\]
Substituting \(v = (\dot{\delta} + DP_0\delta)/H\) from \eqref{eq:rc}:
\[
  \ddot{\delta} + \left(DP_0 + \frac{\zeta}{\tau}\right)\dot{\delta}
  + \frac{DP_0\zeta + HP_0}{\tau}\,\delta = 0.
\]
This is the standard telegraph equation with \(2\gamma = DP_0 + \zeta/\tau\) and \(\omega_0^2 = (DP_0\zeta + HP_0)/\tau\).

\section{Derivation of PME Exponents}

Under \(\delta = \rho_{\max} - \rho\), the ED operator (with \(P_0 = 0\), \(H = 0\)) becomes:
\[
  \partial_t\delta = DM_0\,\nabla\!\cdot\!(\delta^\beta\nabla\delta)
  = \frac{DM_0}{\beta+1}\nabla^2(\delta^{\beta+1}).
\]
This is \(\partial_t u = D_{\mathrm{pme}}\nabla^2(u^m)\) with \(m = \beta + 1\) and \(D_{\mathrm{pme}} = DM_0/(\beta+1)\).

The Barenblatt solution in \(d\) dimensions has front radius exponent \(\alpha_R = 1/(d(m-1)+2)\) and central density exponent \(\alpha_\rho = -d/(d(m-1)+2)\).

\section{Horizon Detection Algorithm}

At each snapshot, compute \(M(\rho(x)) = M_0(\rho_{\max} - \rho(x))^\beta\) at every grid point. The horizon is the connected region where \(M < M_{\mathrm{crit}}\) (default \(M_{\mathrm{crit}} = 0.01\)). The horizon radius \(R_H\) is the maximum distance from the domain centre to any point in the horizon region. If no point satisfies \(M < M_{\mathrm{crit}}\), then \(R_H = 0\).

\section{Temporal Tension}

Two peaks at positions \(x_1\), \(x_2\) with \(s = x_2 - x_1 > 0\). Each peak is tracked by computing the centre of mass of the density excess \(\rho - \rho^*\) in the left (right) half of the domain. The drift rate is \(\dot{s} = ds/dt\), estimated by linear regression on \(s(t)\). The \(v\)-drift correlation is the Pearson correlation between \(ds/dt\) and \(v(t)\) after subtracting means.

The baseline repulsion at \(H = 0\) arises because the degenerate mobility makes diffusion slower where the tails overlap (higher density between the peaks), causing each peak to spread asymmetrically---faster outward, slower inward. This shifts the centre of mass outward, producing an effective repulsion that decays with~\(d\).

%% ============================================================
\begin{thebibliography}{9}

\bibitem{ed-phys}
A.~Proxmire, ``ED-Phys-01 through ED-Phys-40,'' Event Density Physics Series (2025--2026).

\bibitem{ed-sim}
A.~Proxmire, ``ED-SIM-02: An Architectural Lab for Entropic Dynamics,'' Software package (2026).

\bibitem{ed-experiments}
A.~Proxmire, ``ED-PHYS-01 through ED-PHYS-10,'' Event Density Physics Experiments (2026).

\bibitem{ed-bao}
A.~Proxmire, ``Event Density as an Ontology: A BAO-like Correlation Feature from a Minimal PDE'' (2026).

\bibitem{eisenstein}
D.~J.~Eisenstein \emph{et al.}, ``Detection of the baryon acoustic peak in the large-scale correlation function of SDSS luminous red galaxies,'' \emph{Astrophys.\ J.}\ \textbf{633}, 560--574 (2005).

\bibitem{vazquez}
J.~L.~Vazquez, \emph{The Porous Medium Equation: Mathematical Theory}, Oxford University Press (2007).

\bibitem{jko}
R.~Jordan, D.~Kinderlehrer, and F.~Otto, ``The variational formulation of the Fokker--Planck equation,'' \emph{SIAM J.\ Math.\ Anal.}\ \textbf{29}, 1--17 (1998).

\end{thebibliography}

\end{document}
